\chapter{Single-Species Transport}
\def\shorttitle{Single-Species Transport} %

\section{MT3DMS Exercise 1, Part 1: Single-Species 2-D Transport Simulation}
The main purpose of this first exercise is to learn the use of MODFLOW and MT3DMS 
for the analysis of a 2-D contaminant transport problem under both 
natural-gradient and under manipulated hydraulic conditions. 
The effect of different transport solution techniques on 
plume migration will be illustrated.

\subsection{Background: Groundwater Flow}
A simple synthetic aquifer system is considered.
A sketch of the problem is shown in Figure \ref{fig:single_spec_fig1}. 
The aquifer is unconfined and drained by a river. We consider that the flow is mainly perpendicular
to the river. In the numerical model a fixed head boundary can be assumed 
for the western boundary of the domain. The southern and northern boundaries 
can for simplicity be considered as no-flow boundaries assuming that the flow is mainly 
oriented from the West towards the East. 
The model domain has an extension of 250 m in East-West direction and 120 m in North-South direction.
A hydraulic head of 11 m represents the situation at the western boundary. 
A fixed hydraulic head of 10.0 m represents the river water level, assuming that 
the aquifer and the river have a good hydraulic connection. 
A uniform background horizontal hydraulic conductivity of 10 m/d is assumed. 
A zone with a a slightly higher hydraulic conductivity prevails in a zone of old 
river channel deposits. A drinking water extraction 
well is located in the South-East corner of the model domain.


\subsection{Background: Solute Transport}

\begin{figure}[tb]
\centering
\includegraphics[width=1.0\textwidth, angle = 0]{mt3dms_ex1_map}
%\includegraphics[height=0.3]{mt3dms_ex1_map.eps}
%\includegraphics{mt3dms_ex1_map.eps}
\caption[Schematic illustration of the contaminated site]{Schematic illustration of the contaminated site}
\label{fig:single_spec_fig1}
\end{figure}


\begin{figure}[tb]
\centering
\includegraphics[width=1.0\textwidth, angle = 0]{Ex1_Qgis_fg3.eps}
%\includegraphics[height=0.3]{mt3dms_ex1_map.eps}
%\includegraphics{mt3dms_ex1_map.eps}
\caption[Test]{Test}
\label{fig:Test}
\end{figure}

The contaminant source is represented by a small region in the North-West of the domain. 
The concentration in the source zone is 10 mg/L throughout the whole zone,
as indicated in Figure \ref{fig:single_spec_fig1}. 
Clean water enters the upstream boundary, eventually draining to the river 
under natural flow conditions. The boundary conditions along the northern 
and southern model boundaries are implicitly defined as zero mass flux conditions. 
All relevant physical transport parameters are listed in Table \ref{single_spec_table_1}.
They are assumed to be uniform throughout the whole model.
In the exercise we will use the model to determine the maximum allowable 
pumping rate in the extraction well such that the contaminant 
concentrations in that well will remain below 10 \micro g/L.
We will test various model scenarios and assess the affect of the model discretisation
on the model results and of the selection of the solution technique for advectice transport.

\begin{table}[hb]
\caption{Flow and solute transport parameters used for MT3DMS Exercise 1.}
\begin{center}
\begin{tabular}{lc} 
\toprule
Parameter                                            &         Value \\
\midrule
Flow simulation type                                 &        steady state \\
Simulation time (days)                               &  > 800 (until stable) \\
Model \  extent\ column\ direction\ $ (m)$           &  250            \\
Model \  extent\ row\ direction\ $ (m)  $            &  120            \\
Model \  thickness\ $(m)$                            &   15            \\
Model \ top\ elevation\ $(m)$                        &   15            \\
Horizontal hydraulic conductivity (background)\ $(m/day)$         &   10     \\
Horizontal hydraulic conductivity (old channel)\ $(m/day)$         &   25     \\  
Porosity $\theta$                                   &           0.3  \\
Piezometric head upstream boundary\ $(m) $           &  11        \\ 
Piezometric head downstream boundary\ $(m) $         &  10            \\ 
Longitudinal dispersivity $\alpha_L$                &           1.0 m \\
Horizontal transverse dispersivity $\alpha_{TH}$    &           0.1 m \\
Effective diffusion coefficient D$^{\star}$          &             0  \\
\bottomrule

\end{tabular}
\label{single_spec_table_1}
\end{center}
\end{table}

\subsection{Setting up the flow model}

To set up the flow model proceed as follows:

\begin{itemize}

\item In ORTI, create a new model and save it in a new and separate directory (folder).

\item Select {\color{blue}Model Parameters} $\rightarrow$ {\color{blue}1.Model} 
$\rightarrow$ {\color{blue}Model} to specify the model dimension {\color{blue}('2D')} and type {\color{blue}('Free')}

\item Select {\color{blue}Model Parameters} $\rightarrow$ {\color{blue}1.Model} 
$\rightarrow$ {\color{blue}Grid} to specify
the model grid as described in Table \ref{single_spec_table_1}. For the first simulations 
select a spacing of 5 m in x direction and 2.5 m in y direction

\item Set the Top of the aquifer in {\color{blue}Spatial Variables} $\rightarrow$ {\color{blue}MODFLOW} 
$\rightarrow$ {\color{blue}dis.6}, and set input for {\color{blue}one\_value} to 15 m. 
Do not forget to click on the {\color{blue}OK} button. 
The bottom is already set to the default value of 0 and does not need to be modified.

%\item In order to correctly position the spatial objects involved in this model, 
%import a simple map by going to {\color{blue}Model Parameters} $\rightarrow$ 
%{\color{blue}1.Model} 
%\rightarrow$ 
%{\color{blue}Map.} 
%and select the file 
%{\color{blue}mt3dms\_ex1.png}.
%Note, that only png files can be imported as background maps.
%The map must have the same size as the domain and will be positioned directly with 
%its corner at the model boundaries. The map can either be active or hidden under
%the {\color{blue}Results Panel} $\rightarrow$ {\color{blue}1.Model} $\rightarrow$ {\color{blue}Map}.

\item Select {\color{blue}Model Parameters} $\rightarrow$ {\color{blue}Time} $\rightarrow$ 
and set the  {\color{blue} Total Simulation Time} to 800 days. 
Then define the {\color{blue}Step Size} to 25 days. Leave the {\color{blue}Step Mode} as Linear. 

%\item Set the boundary conditions for heads by 
%selecting {\color{blue}Spatial Variables} $\rightarrow$ {\color{blue}Modflow} $\rightarrow$ 
%{\color{blue}BAS6} $\rightarrow$ {\color{blue}bas.3 Boundary Conditions} by using a (background)
%IBOUND (Modflow) value of 1 to define active model grid cells. 
%Use the {\color{blue}Zone Tool} in the {\color{blue}Add Zone} panel to add zones 
%of fixed head BC. For this purpose select the line tool, go into the domain and click
%the first point around coordinates (0.5, 119) and the second at (0.5, 0.5). 
%This will draw a line at the inflow boundary of the domain.
%When the second point is selected (clicked), a dialog appears that allows to set the 
%value of the zone. Set the value to -1 (fixed head BC). 
%Leave the value for {\color{blue}Media} at the default value of 1 (active grid cell).

\item Set the boundary conditions for heads by 
selecting {\color{blue}Spatial Variables} $\rightarrow$ {\color{blue}Modflow} $\rightarrow$ 
{\color{blue}bas.3}, type select zone and a layer called Lbas.3 will appear in the list of Qgis layers, 
then create the zone with Qgis tools

%in spatial variables/Modflow/bas.3 type select zone, a layer called Lbas.3 will appear int he list of Qgis layers, %then create the zone with Qgis tools.

%\item In the same dialog, create another line to represent the fixed head boundary for the river. 
%For this purpose select the {\color{blue}Polygon Tool} 
%and simply click along the river. 
%To end the drawing of the zone right-click on the last point. 
%The value to be entered is again -1.

\item Use the same approach again (using Qgis tools) to create a boundary condition along the river

\item To allocate the correct head values at the fixed head boundaries, 
use {\color{blue}Spatial Variables} $\rightarrow$ {\color{blue}MODFLOW} $\rightarrow$ 
{\color{blue}BAS6}
$\rightarrow$ {\color{blue}bas.5 Initial Heads}. 
Define a value of 10 in the dialog box.
In order to set a value in Qgis in this type of dialog do not forget to type 
ENTER after the value was entered in order to save the value.
This value will be used throughout the domain as a starting head for the computations.
Note, that this value will not change during the iterations at all location in the model grid for which 
the boundary was previously defined to be a fixed head boundary condition (i.e., the value was set to -1).
In order to specify the fixed head values, draw a line at the same place as the inflow BC 
and attribute a value of 11 m. As long as the model grid is relatively 
coarse it is not necessary to put the boundary condition (BC) definition
and the heads definition at exactly the same place. 
However if a more refined grid is used, this will be necessary. 
It is possible, while editing a layer, to select a specific zone and when inspecting its attributes, 
to copy the displayed coordinates and then paste them later when defining the another zone of similar geometry.

%In this case it is possible to select one {\color{blue}Zone}, open the {\color{blue}Zone} dialog, 
%and the copy the coordinates and reuse/paste them into new {\color{blue}Zone} 
%anywhere in the model.

%\item To define the horizontal hydraulic conductivity value 
%select {\color{blue}Spatial Variables} 
%$\rightarrow$ {\color{blue}MODFLOW} $\rightarrow$ {\color{blue}LPF} $\rightarrow$ 
%{\color{blue}lpf.8 Horizontal Hydraulic Conductivity} and set the value to 10 m/d.

\item Again, create a Qgis layer by selecting type zone and then edit this Qgis layer to create a polygon

\item To define the horizontal hydraulic conductivity value in the old river channel zone
select {\color{blue}Spatial Variables} $\rightarrow$ {\color{blue}MODFLOW} $\rightarrow$ 
{\color{blue}LPF} $\rightarrow$ {\color{blue}lpf.8 Horizontal Hydraulic Conductivity} 
and use the {\color{blue}Polygon Tool} to define the outline of the zone of higher 
hydraulic conductivity and set the value to 25 m/d. Another reminder that in the Qgis dialog, 
don't forget to press ENTER for a specified value to be saved.
 

\item To define a (uniform) porosity value, select {\color{blue}Spatial Variables} 
$\rightarrow$ {\color{blue}MT3DMS} $\rightarrow$ {\color{blue}btn.11} 
and set the value to 0.3. 

\item For the inital model runs do not (yet) specify the well.

\item Save the model such that your recent changes will not be lost.

\end{itemize}



\subsection{Running the flow model}

You have now completed to set up the flow model and you are ready to run MODFLOW.
This step will provide MT3DMS with the flow field that is used to compute 
advective-dispersive transport.

\begin{itemize}

\item Select {\color{blue}Model Parameters} $\rightarrow$ {\color{blue}2.Flow} $\rightarrow$ and press 
the {\color{blue}Write Button}.

\item Select {\color{blue}Model Parameters} $\rightarrow$ {\color{blue}2.Flow} $\rightarrow$ and press
the {\color{blue}Run Button}.

\item When the calculation has finished, a window appears, which shows the last lines of the 
main MODFLOW output file. Verify that the last line shows the total simulation time (here 800 days).  
If this is not the case then there was an error in the run and the model input needs to be reassessed and debugged.


\end{itemize}

\subsection{Visualising hydraulic heads}

\begin{itemize}

\item Visualise your flow simulation results by selecting 
{\color{blue}Results} $\rightarrow$ {\color{blue}2.Flow} $\rightarrow$ {\color{blue}Head}. 
The {\color{blue}Time Step} that is to be visualised can be changed in {\color{blue}Results} 
$\rightarrow$ {\color{blue}1.Model} $\rightarrow$ {\color{blue}Time}.

\end{itemize}

The contours should show a steady-state head distribution similar to that 
shown in Figure \ref{fig:single_spec_ex1_heads}.

\begin{figure}[tb]
\centering
\includegraphics[width=1.0\textwidth, angle = 0]{mt3dms_ex1_heads.eps}

\caption[Hydraulic head distribution under natural-gradient conditions.]{Hydraulic head distribution under natural-gradient conditions.}
\label{fig:single_spec_ex1_heads}
\end{figure}


\subsection{Setting up the transport model}

\begin{itemize}

\item Set the boundary conditions for concentrations by selecting 
{\color{blue}Spatial Variables} $\rightarrow$ {\color{blue}MT3DMS} 
$\rightarrow$ {\color{blue}btn.12 Boundary Conditions} (this is the ICBUND value). 
The value must be equal to 1 (= active cell) throughout the domain. In addition 
a {\color{blue}Zone} must be created (a polygon) to approximately replicate 
the source and the value for this {\color{blue}Zone} must be set to -1 
to specify a constant concentration boundary condition.

\item Define the initial concentrations by selecting 
{\color{blue}Spatial Variables} $\rightarrow$ {\color{blue}MT3DMS} 
$\rightarrow$ {\color{blue}btn.13 Concentrations}. 
Use the same coordinates as defined previously to allocate 
an initial concentration of 10 mg/L to this {\color{blue}Zone}.

\item Define the advection package by selecting 
{\color{blue}Model Parameters} $\rightarrow$ {\color{blue}3.Transport} $\rightarrow$ 
{\color{blue}Model Parameters}. In the dialog box choose {\color{blue}ADV} and  
{\color{blue}adv.1} then {\color{blue}MIXELM: Method}. 
Select {\color{blue}TVD} for the initial transport model runs.

\item Define the dispersivity to be used by selecting (in the same dialog), 
the {\color{blue}DSP} $\rightarrow$ {\color{blue}dsp.2} option and 
enter a value of 1 for the longitudinal 
dispersivity. Enter also 0.1 for TRPT under {\color{blue}DSP} $\rightarrow$ {\color{blue}dsp.3}. 

\item Run the transport model by selecting 
{\color{blue}Model Parameters} $\rightarrow$ 
{\color{blue}3.Transport} $\rightarrow$ and pressing the {\color{blue}Write Button}. 
This will write all required {\color{blue}MT3DMS} 
input files. Then press the {\color{blue}Run Button} to start the transport simulation.

\item A MsDos window will appear and you will see rapidly streaming numbers 
until the simulation is finished. After the end of the simulation
you should see the message $END$ $OF$ $MODEL$ $OUTPUT$. 

\end{itemize}

\subsection{Breakthrough curves and contour plots}

To visualise the breakthrough curve for the contaminant at a specific
location proceed as follows:

\begin{itemize}

\item Visualize the concentration contours under
{\color{blue}Results} $\rightarrow$ {\color{blue}3.Transport} $\rightarrow$ {\color{blue}Tracer}. 
You can change the {\color{blue}Time Step} to visualise the results for other simulation times. 
Also, after selecting a time it is possible to use the mouse wheel to changes times.

\item Define an observation well by going to 
{\color{blue}Spatial Variables} $\rightarrow$ {\color{blue}OBSERVATION} 
$\rightarrow$ {\color{blue}obs.1 Observations } and 
use the {\color{blue}Zone Tool} to draw a point at the desired location.

\item Visualise the breakthrough curve by selecting 
{\color{blue}Results} $\rightarrow$ {\color{blue}5.Observation} $\rightarrow$ {\color{blue}Breakthrough}, 
then select the name of the {\color{blue}Observation} point you just created. 
A dialog will appear. Select {\color{blue}Value} and the breakthrough curve for this 
location will appear. The data used for the curve can be exported and used in a spreadsheet
or other postprocessing tools.

\end{itemize}

\begin{figure}[tb]
\centering
\includegraphics[width=0.7\textwidth, angle = 0]{mt3dms_ex1_btc.eps}
\caption[Concentration breakthrough curve]{Concentration breakthrough curve}
\label{fig:single_spec_ex1_btc}
\end{figure}

\begin{figure}[tb]
\centering
\includegraphics[width=1.0\textwidth, angle = 0]{mt3dms_ex1_plume}
%\includegraphics[height=0.9\textwidth, angle = 270]{mt3dms_ex1_2dcont.eps}
\caption[Concentration contours after a 800 day simulation time under 
pumping conditions]{Concentration contours after a 800 day simulation time under 
pumping conditions under pumping conditions}
\label{fig:mt3dms_ex1_2dcont}
\end{figure}

\subsection{Impact pumping well on contaminant plume behaviour}

\begin{itemize}

\item At the location of the well (indicated in the map), 
add a well in the model by going to 
{\color{blue}Spatial Variables} $\rightarrow$ {\color{blue}MODFLOW} $\rightarrow$ 
{\color{blue}wel.1} and create a (point) {\color{blue}Zone}. 
Select a first estimate for the discharge rate of the well. 
Note that negative rates are used for extraction (discharge)
and positive values are used for recharge rate. 
The unit used in this excercise is $m^3/d$.

\item Rewrite the {\color{blue}MODFLOW} input files (to do so press the {\color{blue}Write Button}
and then rerun the flow model.

\item Also rewrite the {\color{blue}MT3DMS} input files and rerun the transport model.

\item To visualise the simulated concentrations at the well an observation point must be
set at the location of the extraction well.

\item Run the model multiple times to determine the maximum pumping rate for which 
the contaminant concentration remains under 10$\micro$g/L.  

\end{itemize}

\subsection{Effect of Solution Techniques}

Now let us first examine the effect of the employed transport solution technique 
on the accuracy of simulation results.

\begin{itemize}
 
\item Export the breakthrough curve obtained at the observation point as obtained when using
the TVD advection option to an ASCII file. 

\item Then change the advection solution option to Upstream Finite-Difference, rerun the model 
and also to MOC, respectively. Export the calculated concentrations to new ASCII files. 
Note that dispersion and sink/source terms are always solved using the standard 
Finite-Difference method and the GCG implicit matrix solver.
Use Excel or other tools to graph the three breakthrough curves.

\end{itemize} 

The TVD and MOC solutions are nearly identical, as both solutions are only 
minimally affected by numerical dispersion. The Upstream Finite-Difference solution, 
on the other hand, shows a significant amount of numerical dispersion.
In addition, check the mass balance errors for each solution. TVD and FD should have
close to zero mass balance errors, while MOC has a large mass balance error at the
beginning. However it should diminish quickly, which is typical for 
particle-tracking based solution techniques.

\subsection{Varying the grid resolution}

With iPHT3D it is possible to easily vary the model discretisation
without the need to change any other model parameters.
To test this effect on the breakthrough curve otained with the Finite-Difference solution technique, 
go to {\color{blue}Model Parameters} $\rightarrow$ {\color{blue}1.Model} 
$\rightarrow$ {\color{blue}Domain} and change the grid cell size.	 

\subsection{Simulating simple contamination remediation methods}
%	 
Go back to the case in which no drinking water extraction takes place and use the model to simulate some basic contaminant remediation and/or attenuation scenarios.
\begin{enumerate}
	\item \textbf{Pump-and-treat}: Illustrate how the contaminant plume could be
      intersected by adding a pump and treat system at >= 80 m downstream of the source zone. Determine the
      pumping rate that would be required to prevent contaminant discharge to the river. Consider 
	  installing multiple wells.
	  %
	\item \textbf{Encapsulate the source zone with a low-permeability barrier}: Illustrate
	  the impact a physical barrier constructed of low-permeability material downstream of the contaminant
	  plume source and/or encapsulate the source zone with low permeability material. Illustrate the impact
	  on both the groundwater flow and on contaminant transport behavior. Experiment with the design of 
	  the barrier and create some figures that document the efficiency of your designs on preventing, 
	  delaying or minimizing contaminant discharge to the river.
	  %	  
	\item \textbf{Natural attenuation scenario}: Add a first-order biodegradation reaction as a chemical
	  reaction into the contaminant transport simulation. Determine the first-order reaction rate constant
	  (and the corresponding half-life) that is required such that the contaminant concentration $C_{dis}$ in
	  the groundwater discharging to the river does not exceed 0.001 $C_{source}$ (ie the concentration you
	  defined at the contaminant source).

\end{enumerate}